% Pedagogical "basics -> advanced" sidebar.  Sits opposite Fig.
% fig:focus-explainer so a reader can read text + visual together.
% Designed to fill ~half a column.

\paragraph{Reading guide -- basic to advanced.}
The FOCUS assigner can be understood at three depths; we encourage the
reader to skim Fig.~\ref{fig:focus-explainer} in parallel.

\textbf{Basic.}\;The assigner is a function that turns a bag of OCR
words into four field values. A receipt has $N\!\approx\!24$ words; the
assigner picks the best word (or word span) for each of
\textsf{company}, \textsf{date}, \textsf{address}, \textsf{total}.

\textbf{Intermediate.}\;Four learned vectors -- one per field -- play
the role of \emph{queries}. Each OCR token is a \emph{key}. We compute
a similarity score between every query and every key, normalise per
field with softmax, and read off a probability distribution
$\alpha_{f,i}$ over which token belongs to field $f$
(Fig.~\ref{fig:focus-internals}).

\textbf{Advanced.}\;Training optimises the \emph{multi-instance
positive-mass NLL}: a field's loss is the negative log of the total
attention mass placed on \emph{any} ground-truth token for that field
(Algorithm~\ref{alg:focus}). This is permutation-invariant over the
positive set $P_f$ -- crucial for multi-token fields like \textsf{address}
where the GT spans 5--7 tokens. A 6-d handcrafted prior
($\phi_i$: \texttt{is\_digit}, \texttt{is\_currency}, \texttt{date\_pattern},
\texttt{y\_pos}, \texttt{x\_pos}, \texttt{height}) gates the keys so the
optimiser does not have to rediscover that ``starts with RM'' implies
\textsf{total}. The hierarchical L1$\to$L2 head
(Fig.~\ref{fig:focus-hier}) cuts the per-field negative set in half,
which calibrates marginals (ECE drops from $0.094 \to 0.041$).
